\documentclass[11pt, a4paper]{article}

\usepackage[utf8]{inputenc}
\usepackage[T1]{fontenc}
\usepackage{amsmath, amssymb, amsthm, mathtools}
\usepackage{bm}
\usepackage{booktabs}
\usepackage{graphicx}
\usepackage{hyperref}
\usepackage[margin=2.5cm]{geometry}
\usepackage{enumitem}
\usepackage{xcolor}
\usepackage{tcolorbox}
\usepackage{float}

% Algorithm box (no external package needed)
\newtcolorbox{algorithmbox}[1][]{
    colback=white,
    colframe=black!70,
    fonttitle=\bfseries,
    #1
}

% Theorem environments
\theoremstyle{definition}
\newtheorem{definition}{Definition}[section]
\newtheorem{example}{Example}[section]
\theoremstyle{plain}
\newtheorem{theorem}{Theorem}[section]
\newtheorem{proposition}{Proposition}[section]
\newtheorem{corollary}{Corollary}[section]
\newtheorem{lemma}{Lemma}[section]
\theoremstyle{remark}
\newtheorem{remark}{Remark}[section]

% Colored boxes
\newtcolorbox{keyresult}[1][]{
    colback=blue!5!white,
    colframe=blue!60!black,
    fonttitle=\bfseries,
    title={Key Insight},
    #1
}
\newtcolorbox{warningbox}[1][]{
    colback=red!5!white,
    colframe=red!60!black,
    fonttitle=\bfseries,
    title={Warning},
    #1
}
\newtcolorbox{refbox}[1][]{
    colback=gray!5!white,
    colframe=gray!50!black,
    fonttitle=\bfseries,
    title={Go Deeper},
    #1
}
\newtcolorbox{keyidea}[1][]{
    colback=green!5!white,
    colframe=green!50!black,
    fonttitle=\bfseries,
    title={Core Idea},
    #1
}
\newtcolorbox{prereq}[1][]{
    colback=orange!5!white,
    colframe=orange!60!black,
    fonttitle=\bfseries,
    title={Read Before Coding},
    #1
}
\newtcolorbox{objective}[1][]{
    colback=blue!5!white,
    colframe=blue!60!black,
    fonttitle=\bfseries,
    title={What You Will Build},
    #1
}

% Custom commands
\newcommand{\VaR}{\operatorname{VaR}}
\newcommand{\CVaR}{\operatorname{CVaR}}
\newcommand{\ES}{\operatorname{ES}}
\newcommand{\EV}{\mathbb{E}}
\newcommand{\PV}{\mathbb{P}}
\newcommand{\RV}{\mathbb{R}}
\newcommand{\ind}{\boldsymbol{1}}

\title{%
    \textbf{Climate Risk Scenario Analysis} \\[0.5em]
    \large Transition Risk, Physical Risk, and TCFD-Aligned Reporting
}
\author{Andr\'es Gonz\'alez Ortega}
\date{March 2026}

\begin{document}
\maketitle

\begin{abstract}
An applied treatment of climate risk scenario analysis for financial portfolios,
covering both transition risk and physical risk under the NGFS scenario framework.
Climate change poses a dual threat to financial stability: transition risk from the
repricing of carbon-intensive assets during the shift to a low-carbon economy, and
physical risk from the direct economic damages of rising temperatures, extreme weather
events, and chronic environmental degradation.
This document develops both risk channels from first principles.
We formalize transition risk through carbon cost propagation and stranded asset
valuation; physical risk through the Nordhaus damage function and extreme value theory
for acute climate events.
These are combined into a unified Climate Value-at-Risk metric spanning multiple
NGFS pathways.
A Sobol sensitivity analysis identifies the dominant risk drivers, and TCFD-aligned
reporting metrics --- Weighted Average Carbon Intensity, financed emissions, and
implied temperature rise --- are derived for portfolio-level disclosure.
We assume familiarity with basic portfolio theory, Value-at-Risk, and the extreme
value theory foundations from Project~05.
\end{abstract}

\tableofcontents
\newpage

%% ============================================================================
\section{Why Climate Risk Matters for Finance}
\label{sec:why-climate}
%% ============================================================================

Climate risk is no longer a reputational or ethical concern relegated to ESG
departments.
It is a \emph{financial} risk --- one that regulators, central banks, and
standard-setters now treat as material to portfolio valuation, capital adequacy,
and systemic stability.

\subsection{The Regulatory Shift: From Voluntary to Mandatory}

The Task Force on Climate-related Financial Disclosures (TCFD), established by the
Financial Stability Board in 2017, laid the conceptual groundwork: firms should
disclose how climate change affects their strategy, governance, risk management, and
financial metrics.
The TCFD framework was initially voluntary.
That era is over.

\begin{itemize}
    \item \textbf{ISSB (2023).} The International Sustainability Standards Board
          published IFRS~S2 (\emph{Climate-related Disclosures}), making
          TCFD-aligned reporting a baseline for capital markets globally.
          Jurisdictions including the UK, EU, Japan, and Australia have moved to
          adopt or mandate ISSB standards.
    \item \textbf{ECB (2025).} The European Central Bank's 2025 macroprudential
          stress test required Eurozone banks to model portfolio losses under
          NGFS climate scenarios, integrating transition and physical risk channels
          into supervisory capital assessments for the first time at scale.
    \item \textbf{EU CSRD (2024).} The Corporate Sustainability Reporting Directive
          requires approximately 50{,}000 companies to report under the European
          Sustainability Reporting Standards, which embed double materiality ---
          both how climate affects the company and how the company affects climate.
\end{itemize}

\begin{warningbox}
The transition from voluntary to mandatory disclosure means that climate risk
modeling is no longer optional for financial institutions.
A bank that cannot quantify its portfolio exposure to a \$200/tCO$_2$ carbon price
or a 3$^\circ$C warming scenario faces both regulatory penalties and market
repricing once peers begin disclosing.
\end{warningbox}

\subsection{Double Materiality}

The concept of \textbf{double materiality} is central to the European framework
and increasingly influential elsewhere:

\begin{definition}[Double materiality]\label{def:double-materiality}
A risk is \textbf{doubly material} if it operates in both directions:
\begin{enumerate}
    \item \textbf{Outside-in (financial materiality):} Climate change affects the
          value of financial assets --- through carbon pricing, physical damages,
          stranded assets, or liability risks.
    \item \textbf{Inside-out (impact materiality):} The portfolio's financed
          activities contribute to climate change --- through greenhouse gas
          emissions, deforestation, or resource depletion --- creating systemic
          risk that feeds back into financial outcomes.
\end{enumerate}
\end{definition}

A coal-heavy portfolio is exposed from the outside-in (carbon prices reduce
earnings) \emph{and} contributes from the inside-out (financed emissions
accelerate warming, increasing physical risk for all market participants).

\subsection{The Scale of Exposure}

The financial stakes are large:
\begin{itemize}
    \item \textbf{Stranded assets.} An estimated \$2.5 trillion in fossil fuel
          reserves become economically unviable if carbon prices rise to levels
          consistent with the Paris Agreement.
          These reserves are currently on corporate balance sheets, underpinning
          equity valuations and debt covenants.
    \item \textbf{Physical damages.} Under a 3$^\circ$C+ warming pathway, the
          Nordhaus damage function implies cumulative GDP losses of 4\% or more ---
          concentrated in agriculture, insurance, real estate, and infrastructure.
    \item \textbf{Transition costs.} An orderly transition to net zero by 2050
          requires annual clean-energy investment of approximately \$4 trillion
          (IEA, 2023), implying massive capital reallocation away from
          carbon-intensive sectors.
\end{itemize}

\begin{keyidea}
Climate risk is not a single number.
It is a \emph{scenario-dependent} distribution of losses that varies with the
speed and ambition of the low-carbon transition.
An orderly transition imposes moderate, predictable transition costs but limits
physical damages.
A delayed or absent transition avoids near-term transition costs but exposes
portfolios to severe, non-linear physical damages.
The challenge for risk managers is to evaluate portfolios under \emph{all}
plausible pathways.
\end{keyidea}


%% ============================================================================
\section{The NGFS Scenario Framework}
\label{sec:ngfs}
%% ============================================================================

The Network for Greening the Financial System (NGFS) provides the reference
scenario set used by most central banks and supervisory authorities.
Phase~V (2024) defines six pathways organized along two axes: the speed and
coordination of transition policy (orderly versus disorderly) and the level of
physical risk (determined by the resulting warming outcome).

\subsection{The Six Pathways}

\begin{table}[H]
\centering
\caption{NGFS Phase~V scenario matrix. Temperature outcomes refer to global
mean surface temperature increase by 2100 relative to pre-industrial levels.
Shadow carbon prices are illustrative 2035 values in USD/tCO$_2$.}
\label{tab:ngfs-scenarios}
\begin{tabular}{@{}p{3.2cm}p{2.0cm}p{1.8cm}p{1.8cm}p{3.5cm}@{}}
\toprule
\textbf{Scenario} & \textbf{Transition Risk} & \textbf{Warming ($^\circ$C)} &
\textbf{Carbon Price (\$/tCO$_2$)} & \textbf{Character} \\
\midrule
Net Zero 2050       & Orderly    & 1.5 & $\sim$300 & Immediate, coordinated \\[4pt]
Below 2$^\circ$C    & Orderly    & 1.7 & $\sim$200 & Gradual strengthening \\[4pt]
Low Demand          & Orderly    & 1.6 & $\sim$150 & Demand-side shift \\[4pt]
Delayed Transition  & Disorderly & 1.8 & $\sim$350$^*$ & Late, abrupt policy \\[4pt]
NDCs                & Insufficient & 2.6 & $\sim$50 & Current pledges only \\[4pt]
Current Policies    & Minimal    & 3.0+ & $\sim$10  & No new policy \\
\bottomrule
\end{tabular}

\smallskip
\footnotesize{$^*$Delayed Transition carbon price is low until $\sim$2030, then
spikes sharply --- the disorderly character implies a sudden repricing event.}
\end{table}

\begin{keyresult}
The six scenarios span a fundamental trade-off: \textbf{transition risk and
physical risk are inversely correlated across pathways}.
Orderly scenarios have high carbon prices (transition risk) but limit warming
to 1.5--1.7$^\circ$C (low physical risk).
Current Policies has negligible transition risk but exposes portfolios to
3$^\circ$C+ warming (high physical risk).
The Delayed Transition scenario is particularly dangerous: it combines high
transition risk (sudden policy action) with moderate physical risk (1.8$^\circ$C
warming already locked in).
\end{keyresult}

\subsection{Shadow Carbon Price Trajectories}

The \textbf{shadow carbon price} is the implicit cost of carbon emissions required
to achieve a given temperature target.
Under the Net Zero 2050 pathway:

\begin{equation}\label{eq:carbon-price-path}
    p_{\text{CO}_2}(t) \approx
    \begin{cases}
        50 + 50 \cdot (t - 2025) & \text{for } 2025 \leq t \leq 2030, \\
        300 + 20 \cdot (t - 2035) & \text{for } t > 2035,
    \end{cases}
\end{equation}
reaching approximately \$300/tCO$_2$ by 2035 and continuing to rise thereafter.
Under the Delayed Transition pathway, the carbon price remains below \$30/tCO$_2$
until approximately 2030, then rises sharply to \$350/tCO$_2$ by 2035 --- a
price shock that propagates through energy markets, credit spreads, and equity
valuations simultaneously.

\subsection{Short-Term Scenarios (May 2025)}

NGFS introduced short-term scenarios with a 2025--2030 horizon, reflecting the
recognition that long-run (2050--2100) pathways, while important for strategic
planning, do not capture the near-term financial risks that drive supervisory
stress testing and portfolio risk management.
These short-term scenarios incorporate:
\begin{itemize}
    \item Specific policy commitments (e.g., EU CBAM, US IRA provisions).
    \item Near-term commodity price dynamics and energy supply constraints.
    \item Technology readiness levels for key abatement options.
\end{itemize}

\begin{refbox}
The NGFS scenario portal (\texttt{https://www.ngfs.net/ngfs-scenarios-portal})
provides downloadable variable paths for GDP, carbon prices, energy mix,
temperature, and sectoral output for all six pathways.
These are the direct inputs to the Climate VaR computation in
Section~\ref{sec:climate-var}.
\end{refbox}


%% ============================================================================
\section{Transition Risk Modeling}
\label{sec:transition-risk}
%% ============================================================================

Transition risk arises when the shift to a low-carbon economy changes the
relative costs, revenues, and asset values of firms and sectors.
The primary channel is carbon pricing --- either explicit (carbon taxes, emissions
trading systems) or implicit (regulation, shifts in demand).

\subsection{Carbon Cost Propagation}

The annualized carbon cost for a firm is:

\begin{equation}\label{eq:carbon-cost}
    C_{\text{carbon}} = \underbrace{e_i}_{\text{emission intensity}}
    \times \underbrace{p_{\text{CO}_2}}_{\text{carbon price}}
    \times \underbrace{R_i}_{\text{revenue}},
\end{equation}
where $e_i$ is the firm's emission intensity measured in tonnes of CO$_2$ per
million dollars of revenue (tCO$_2$/\$M), $p_{\text{CO}_2}$ is the carbon price
in dollars per tonne, and $R_i$ is the firm's annual revenue.

\begin{keyidea}
Carbon cost scales \emph{multiplicatively} in two policy-dependent variables:
emission intensity (which firms can reduce through abatement) and carbon price
(which is set by regulation or market forces).
A firm with high emission intensity but low revenue may face less absolute carbon
cost than a moderate-intensity firm with large revenue.
Sector-level analysis must account for both.
\end{keyidea}

\subsection{Sector Emission Intensities}

Emission intensities vary by orders of magnitude across sectors:

\begin{table}[H]
\centering
\caption{Approximate Scope 1+2 emission intensities by sector. Values are
representative and vary by geography and sub-industry.}
\label{tab:emission-intensity}
\begin{tabular}{@{}lrl@{}}
\toprule
\textbf{Sector} & \textbf{Intensity (tCO$_2$/\$M revenue)} &
\textbf{Primary Source} \\
\midrule
Energy (Oil \& Gas)   & $\sim$500 & Extraction, refining, combustion \\
Utilities             & $\sim$300 & Power generation (coal, gas) \\
Materials (Steel, Cement) & $\sim$200 & Industrial processes \\
Transportation        & $\sim$150 & Fuel combustion \\
Agriculture           & $\sim$120 & Livestock, fertilizer, land use \\
Financials            & $\sim$50  & Financed emissions (Scope 3) \\
Technology            & $\sim$10  & Data centers, supply chain \\
\bottomrule
\end{tabular}
\end{table}

\subsection{Equity Repricing Under Carbon Costs}

When a carbon cost is imposed, the firm's earnings decline.
The first-order equity repricing is:

\begin{equation}\label{eq:equity-repricing}
    \Delta V_i = -\frac{C_{\text{carbon},i}}{\text{EBITDA}_i} \times V_i,
\end{equation}
where $V_i$ is the current equity value and $\text{EBITDA}_i$ is earnings
before interest, taxes, depreciation, and amortization.
The ratio $C_{\text{carbon}} / \text{EBITDA}$ measures how much of the firm's
earnings are consumed by carbon costs --- a direct analogue of the interest
coverage ratio for debt.

\begin{definition}[Carbon earnings-at-risk]\label{def:carbon-ear}
The \textbf{carbon earnings-at-risk} for firm $i$ under scenario $s$ is:
\begin{equation}
    \text{CarbonEaR}_{i,s} = \frac{e_i \cdot p_{\text{CO}_2}^{(s)} \cdot R_i}
    {\text{EBITDA}_i}.
\end{equation}
A CarbonEaR above 1.0 implies that carbon costs exceed the firm's entire
operating earnings --- the firm becomes unprofitable under that carbon price.
\end{definition}

\subsection{Stranded Assets}

Fossil fuel reserves are classified as \textbf{stranded} when the cost of
extraction plus the carbon cost exceeds the market price of the commodity:

\begin{equation}\label{eq:stranded}
    \text{Reserve is stranded} \iff
    c_{\text{extract}} + p_{\text{CO}_2} \cdot f_{\text{emission}} > p_{\text{market}},
\end{equation}
where $c_{\text{extract}}$ is the per-barrel (or per-tonne) extraction cost,
$f_{\text{emission}}$ is the emission factor (tCO$_2$ per barrel), and
$p_{\text{market}}$ is the commodity price.

\begin{example}[Stranded reserves under \$200/tCO$_2$]\label{ex:stranded}
Consider an oil producer with the following profile:
\begin{itemize}
    \item Proven reserves: 2 billion barrels of oil equivalent (boe).
    \item Extraction cost distribution: 40\% of reserves at \$30/boe,
          35\% at \$50/boe, 25\% at \$70/boe.
    \item Emission factor: 0.43 tCO$_2$/boe (industry average).
    \item Market oil price: \$75/boe.
\end{itemize}

Under a carbon price of \$200/tCO$_2$, the effective carbon cost per barrel is
$200 \times 0.43 = \$86$/boe.
The total cost (extraction + carbon) for each reserve tranche:
\begin{center}
\begin{tabular}{@{}cccc@{}}
\toprule
\textbf{Tranche} & \textbf{Extraction} & \textbf{Total Cost} &
\textbf{Status} \\
\midrule
40\% at \$30/boe & \$30 & \$30 + \$86 = \$116 & Stranded ($>\$75$) \\
35\% at \$50/boe & \$50 & \$50 + \$86 = \$136 & Stranded ($>\$75$) \\
25\% at \$70/boe & \$70 & \$70 + \$86 = \$156 & Stranded ($>\$75$) \\
\bottomrule
\end{tabular}
\end{center}

At \$200/tCO$_2$, \textbf{100\%} of this producer's reserves are stranded ---
none can be profitably extracted.
Even at a more moderate \$100/tCO$_2$ (carbon cost = \$43/boe), the highest-cost
tranche (\$70 + \$43 = \$113) and mid-cost tranche (\$50 + \$43 = \$93) are
stranded, leaving only 40\% of reserves viable.
The equity write-down from stranding 60\% of reserves at \$100/tCO$_2$ is:
\begin{equation}
    \Delta V = -0.60 \times \text{NPV}(\text{reserves}) \times
    \frac{V_{\text{equity}}}{\text{EV}},
\end{equation}
where EV is enterprise value.
For a firm where reserves represent 70\% of enterprise value, this implies a
$\sim$42\% equity loss.
\end{example}

\begin{warningbox}
Stranded asset risk is highly non-linear.
At low carbon prices, most reserves remain viable.
Above a critical threshold --- which varies by extraction cost profile --- the
fraction of stranded reserves jumps sharply.
This creates \emph{cliff risk}: a small increase in carbon price near the
threshold can trigger a large write-down.
\end{warningbox}


%% ============================================================================
\section{Physical Risk Modeling}
\label{sec:physical-risk}
%% ============================================================================

Physical risk is the direct economic damage from climate change.
It operates through two channels: \textbf{acute} risks (discrete extreme events
whose frequency and severity increase with warming) and \textbf{chronic} risks
(gradual shifts that erode productivity and asset values over decades).

\subsection{The Nordhaus Damage Function}

The workhorse model for aggregate physical damages is the quadratic damage
function introduced by Nordhaus (2017) in the DICE model:

\begin{equation}\label{eq:nordhaus-damage}
    D(T) = \alpha \cdot T^2,
\end{equation}
where $D(T)$ is the fraction of GDP lost at a global mean temperature increase
of $T$ degrees Celsius above pre-industrial levels, and $\alpha$ is a calibration
parameter.
Nordhaus calibrates $\alpha \approx 0.00267$, implying:

\begin{table}[H]
\centering
\caption{GDP loss under the Nordhaus quadratic damage function
($\alpha = 0.00267$).}
\label{tab:damage-function}
\begin{tabular}{@{}ccc@{}}
\toprule
\textbf{Warming ($^\circ$C)} & \textbf{GDP Loss (\%)} & \textbf{NGFS Scenario} \\
\midrule
1.5 & 0.6\% & Net Zero 2050 \\
2.0 & 1.1\% & Below 2$^\circ$C \\
2.5 & 1.7\% & --- \\
3.0 & 2.4\% & Current Policies \\
4.0 & 4.3\% & Current Policies (high sensitivity) \\
\bottomrule
\end{tabular}
\end{table}

\begin{keyidea}
The quadratic functional form is crucial: damages are \textbf{convex} in
temperature.
The marginal damage of each additional degree of warming is \emph{increasing}.
Moving from 2$^\circ$C to 3$^\circ$C is more than twice as costly as moving
from 1$^\circ$C to 2$^\circ$C.
This convexity is the fundamental reason why aggressive mitigation (avoiding
high warming) is economically rational even when discounted at market rates.
\end{keyidea}

\begin{remark}
The Nordhaus damage function has been criticized as too conservative.
Howard and Sterner (2017) survey the empirical literature and find that including
non-market damages (biodiversity loss, health impacts, conflict) and accounting
for persistent growth effects roughly doubles the damage coefficient to
$\alpha \approx 0.005$, implying $\sim$8\% GDP loss at 4$^\circ$C.
Our framework treats $\alpha$ as a parameter to be stress-tested rather than a
fixed constant.
\end{remark}

\subsection{Acute Physical Risk: Extreme Weather Events}

Acute physical risk manifests as discrete, high-severity events: floods,
hurricanes, wildfires, and heat waves.
Climate change increases both the frequency and the severity of these events,
shifting the tail of the loss distribution.

\begin{definition}[Climate-adjusted extreme value distribution]\label{def:climate-evt}
Let $L$ denote the annual maximum loss from a climate-sensitive peril (e.g.,
flood damage).
Under the Generalized Pareto Distribution (GPD) framework from extreme value
theory (see Project~05), exceedances above a threshold $u$ follow:
\begin{equation}\label{eq:gpd-climate}
    \PV(L - u > y \mid L > u) = \left(1 + \xi \frac{y}{\sigma}\right)^{-1/\xi},
    \quad y > 0,
\end{equation}
where $\xi > 0$ is the shape parameter (controlling tail heaviness) and $\sigma > 0$
is the scale parameter.
Under climate change, both parameters are functions of temperature:
\begin{equation}\label{eq:climate-gpd-params}
    \sigma(T) = \sigma_0 \cdot \exp(\beta_\sigma \cdot T), \qquad
    \xi(T) = \xi_0 + \beta_\xi \cdot T,
\end{equation}
reflecting that warming increases both the scale (larger typical extreme losses)
and the shape (heavier tails --- more ``black swan'' events).
\end{definition}

\begin{example}[100-year flood loss under warming]\label{ex:flood-gpd}
Consider flood losses for a coastal property portfolio.
Under current climate ($T = 1.1^\circ$C above pre-industrial):
\begin{itemize}
    \item Threshold: $u = \$50$M (losses exceeding this are ``extreme'').
    \item GPD parameters: $\sigma_0 = 20$, $\xi_0 = 0.15$.
    \item 100-year return level (0.01 annual exceedance probability):
\end{itemize}
\begin{equation}
    L_{100} = u + \frac{\sigma_0}{\xi_0}
    \left[\left(\frac{1}{0.01}\right)^{\xi_0} - 1\right]
    = 50 + \frac{20}{0.15}\left[100^{0.15} - 1\right]
    = 50 + 133.3 \times 0.995 = \$183\text{M}.
\end{equation}

Under 3$^\circ$C warming, with $\beta_\sigma = 0.08$ and $\beta_\xi = 0.01$
per degree:
\begin{itemize}
    \item $\sigma(3.0) = 20 \cdot \exp(0.08 \times 3.0) = 20 \times 1.271 = 25.4$.
    \item $\xi(3.0) = 0.15 + 0.01 \times 3.0 = 0.18$.
\end{itemize}
\begin{equation}
    L_{100}^{3^\circ\text{C}} = 50 + \frac{25.4}{0.18}
    \left[100^{0.18} - 1\right]
    = 50 + 141.1 \times 1.318 = \$236\text{M}.
\end{equation}

The 100-year flood loss increases from \$183M to \$236M --- a \textbf{29\% increase}
--- under 3$^\circ$C warming.
More critically, what was a 100-year event under current climate becomes
approximately a 50-year event, doubling the expected annual frequency.
\end{example}

\subsection{Chronic Physical Risk}

Chronic risks accumulate gradually and affect asset productivity and valuations
over decades:

\begin{itemize}
    \item \textbf{Sea-level rise.} Coastal real estate and infrastructure face
          permanent inundation or increased flood insurance costs.
          Under 3$^\circ$C, global mean sea level rises approximately 0.5--1.0m
          by 2100.
    \item \textbf{Heat stress.} Labor productivity declines as the number of
          days exceeding safe working temperatures increases.
          Outdoor sectors (agriculture, construction, mining) are most exposed.
    \item \textbf{Agricultural yield decline.} Crop yields decline non-linearly
          with temperature beyond optimal growing ranges.
          Each degree above 2$^\circ$C reduces global yields by approximately
          5--10\% for staple crops (IPCC AR6).
    \item \textbf{Water stress.} Changing precipitation patterns affect
          hydropower, agriculture, and industrial water supply.
\end{itemize}

\subsection{Sector Exposure to Physical Risk}

\begin{table}[H]
\centering
\caption{Sector exposure to physical risk channels.
H = High, M = Medium, L = Low.}
\label{tab:physical-exposure}
\begin{tabular}{@{}lcccc@{}}
\toprule
\textbf{Sector} & \textbf{Acute (Floods/Storms)} & \textbf{Sea-Level Rise} &
\textbf{Heat Stress} & \textbf{Yield Decline} \\
\midrule
Agriculture       & H & M & H & H \\
Insurance         & H & H & M & L \\
Real Estate       & H & H & M & L \\
Infrastructure    & H & H & M & L \\
Utilities         & M & M & M & L \\
Transportation    & M & M & L & L \\
Technology        & L & L & L & L \\
\bottomrule
\end{tabular}
\end{table}

\begin{keyresult}
Transition risk and physical risk affect \emph{different} sectors.
Energy, utilities, and materials face the highest transition risk (carbon
intensity).
Agriculture, insurance, and real estate face the highest physical risk
(exposure to weather and temperature).
A well-diversified portfolio is exposed to \emph{both} channels simultaneously,
and the correlation between them depends on the scenario pathway.
\end{keyresult}


%% ============================================================================
\section{Climate Value-at-Risk}
\label{sec:climate-var}
%% ============================================================================

Climate VaR aggregates transition and physical losses into a single portfolio-level
risk metric, evaluated under each NGFS scenario.

\subsection{Definition}

\begin{definition}[Climate Value-at-Risk]\label{def:climate-var}
For a portfolio with $N$ holdings, under NGFS scenario $s$ with time horizon $H$
(typically 10--30 years), the \textbf{Climate VaR} at confidence level $1 - \alpha$
is:
\begin{equation}\label{eq:climate-var}
    \text{ClimateVaR}_\alpha^{(s)} = \inf\left\{l \geq 0 :
    \PV\!\left(\sum_{i=1}^{N} w_i \cdot
    \bigl(L_i^{\text{trans}}(s) + L_i^{\text{phys}}(s)\bigr) \leq l\right)
    \geq 1 - \alpha\right\},
\end{equation}
where $w_i$ is the portfolio weight, $L_i^{\text{trans}}(s)$ is the transition
loss for holding $i$ under scenario $s$, and $L_i^{\text{phys}}(s)$ is the
physical loss.
\end{definition}

\subsection{Loss Components}

The total loss for holding $i$ under scenario $s$ is:
\begin{equation}\label{eq:total-loss}
    L_i^{(s)} = L_i^{\text{trans}}(s) + L_i^{\text{phys}}(s),
\end{equation}
where:

\paragraph{Transition loss.}
\begin{equation}\label{eq:trans-loss}
    L_i^{\text{trans}}(s) = \frac{e_i \cdot p_{\text{CO}_2}^{(s)}(H)
    \cdot R_i}{\text{EBITDA}_i},
\end{equation}
expressed as a fraction of current equity value (i.e., the carbon
earnings-at-risk from Definition~\ref{def:carbon-ear}).

\paragraph{Physical loss.}
\begin{equation}\label{eq:phys-loss}
    L_i^{\text{phys}}(s) = \phi_i \cdot D\!\left(T^{(s)}(H)\right)
    = \phi_i \cdot \alpha \cdot \bigl(T^{(s)}(H)\bigr)^2,
\end{equation}
where $\phi_i \in [0, 1]$ is the physical risk exposure factor for holding $i$
(from Table~\ref{tab:physical-exposure}), $T^{(s)}(H)$ is the cumulative warming
by horizon $H$ under scenario $s$, and $D(\cdot)$ is the Nordhaus damage function.

\subsection{Computation}

For deterministic scenarios (a single pathway per NGFS scenario), the Climate VaR
reduces to a point estimate:
\begin{equation}\label{eq:climate-var-deterministic}
    \text{ClimateVaR}^{(s)} = \sum_{i=1}^{N} w_i \cdot L_i^{(s)}.
\end{equation}

For stochastic scenarios (where carbon prices, temperatures, and damages are
random variables within each pathway), we simulate:

\begin{algorithmbox}[title={Algorithm 1: Stochastic Climate VaR}]
\label{alg:climate-var}
\textbf{Input:} Portfolio weights $\{w_i\}$, firm data $\{e_i, R_i,
\text{EBITDA}_i, \phi_i\}$, NGFS scenario $s$, number of simulations $M$,
confidence level $1 - \alpha$.
\begin{enumerate}
    \item For $m = 1, \ldots, M$:
    \begin{enumerate}[label=(\alph*)]
        \item Draw carbon price: $p_{\text{CO}_2}^{(m)} \sim
              F_{\text{CO}_2}^{(s)}$.
        \item Draw temperature: $T^{(m)} \sim F_T^{(s)}$.
        \item Draw GDP growth path: $g^{(m)} \sim F_g^{(s)}$.
        \item Compute transition loss for each holding:
              $L_i^{\text{trans},(m)} = e_i \cdot p_{\text{CO}_2}^{(m)} \cdot R_i
              / \text{EBITDA}_i$.
        \item Compute physical loss for each holding:
              $L_i^{\text{phys},(m)} = \phi_i \cdot \alpha \cdot (T^{(m)})^2$.
        \item Compute portfolio loss:
              $L^{(m)} = \sum_{i=1}^{N} w_i \cdot (L_i^{\text{trans},(m)} +
              L_i^{\text{phys},(m)})$.
    \end{enumerate}
    \item Sort $\{L^{(1)}, \ldots, L^{(M)}\}$.
    \item Return $\text{ClimateVaR}_\alpha^{(s)} = L^{(\lceil M(1-\alpha) \rceil)}$.
\end{enumerate}
\end{algorithmbox}

\subsection{Comparison with Traditional VaR}

\begin{table}[H]
\centering
\caption{Traditional VaR vs.\ Climate VaR.}
\label{tab:var-comparison}
\begin{tabular}{@{}p{3.5cm}p{5cm}p{5cm}@{}}
\toprule
\textbf{Dimension} & \textbf{Traditional VaR} & \textbf{Climate VaR} \\
\midrule
Time horizon & 1 day -- 10 days & 5 -- 30 years \\
Risk drivers & Market prices, volatility & Carbon price, temperature, policy \\
Data source & Historical returns & NGFS scenario projections \\
Distribution & Estimated from data & Scenario-conditioned simulation \\
Stationarity & Assumed (or windowed) & Explicitly non-stationary \\
Interpretation & ``Loss in normal conditions'' & ``Loss under pathway $s$'' \\
\bottomrule
\end{tabular}
\end{table}

\begin{example}[Portfolio Climate VaR across pathways]\label{ex:climate-var}
Consider a portfolio with three sector exposures:
\begin{itemize}
    \item 30\% Energy ($e = 500$, EBITDA margin = 15\%, $\phi = 0.3$).
    \item 40\% Financials ($e = 50$, EBITDA margin = 35\%, $\phi = 0.4$).
    \item 30\% Technology ($e = 10$, EBITDA margin = 25\%, $\phi = 0.1$).
\end{itemize}

Under the Delayed Transition pathway ($p_{\text{CO}_2} = \$350$, $T = 1.8^\circ$C
by 2050):
\begin{align}
    L_{\text{Energy}}^{\text{trans}} &= \frac{500 \times 350 \times 10^{-6}}{0.15}
    = 1.167 \quad (\text{i.e., 116.7\% of equity}), \\
    L_{\text{Fin}}^{\text{trans}} &= \frac{50 \times 350 \times 10^{-6}}{0.35}
    = 0.050 \quad (5.0\%), \\
    L_{\text{Tech}}^{\text{trans}} &= \frac{10 \times 350 \times 10^{-6}}{0.25}
    = 0.014 \quad (1.4\%).
\end{align}
Physical losses (Nordhaus, $\alpha = 0.00267$, $T = 1.8$):
$D(1.8) = 0.00267 \times 1.8^2 = 0.87\%$.
\begin{align}
    L_{\text{Energy}}^{\text{phys}} &= 0.3 \times 0.0087 = 0.26\%, \\
    L_{\text{Fin}}^{\text{phys}} &= 0.4 \times 0.0087 = 0.35\%, \\
    L_{\text{Tech}}^{\text{phys}} &= 0.1 \times 0.0087 = 0.09\%.
\end{align}
Portfolio Climate VaR:
\begin{equation}
    \text{ClimateVaR} = 0.30 \times 1.170 + 0.40 \times 0.054 + 0.30 \times 0.015
    = 37.7\%.
\end{equation}
Under the delayed transition pathway, this portfolio could lose approximately
38\% of its value over the projection horizon --- driven almost entirely by the
energy sector's transition risk.
\end{example}


%% ============================================================================
\section{Sobol Sensitivity Analysis}
\label{sec:sobol}
%% ============================================================================

Climate VaR depends on multiple uncertain inputs: carbon price, temperature,
GDP growth, discount rates, emission intensities, and physical exposure factors.
\textbf{Sobol sensitivity analysis} provides a rigorous, variance-based
decomposition of which inputs drive the most uncertainty in the output.

\subsection{Variance Decomposition}

Let $Y = f(X_1, X_2, \ldots, X_k)$ be the Climate VaR as a function of $k$
uncertain inputs.
The total variance of $Y$ can be decomposed (under independence of inputs) as:

\begin{equation}\label{eq:anova-decomposition}
    \operatorname{Var}(Y) = \sum_{i=1}^{k} V_i
    + \sum_{i < j} V_{ij} + \cdots + V_{1,2,\ldots,k},
\end{equation}
where $V_i = \operatorname{Var}\bigl(\EV[Y \mid X_i]\bigr)$ is the variance
attributable to $X_i$ alone, $V_{ij}$ captures the interaction between $X_i$ and
$X_j$, and so on.

\subsection{First-Order and Total-Order Indices}

\begin{definition}[Sobol indices]\label{def:sobol}
The \textbf{first-order Sobol index} for input $X_i$ is:
\begin{equation}\label{eq:sobol-first}
    S_i = \frac{V_i}{\operatorname{Var}(Y)}
    = \frac{\operatorname{Var}\bigl(\EV[Y \mid X_i]\bigr)}{\operatorname{Var}(Y)}.
\end{equation}
It measures the fraction of output variance attributable to $X_i$ \emph{alone},
averaging over all other inputs.

The \textbf{total-order Sobol index} for input $X_i$ is:
\begin{equation}\label{eq:sobol-total}
    S_{T_i} = 1 - \frac{\operatorname{Var}\bigl(\EV[Y \mid X_{\sim i}]\bigr)}
    {\operatorname{Var}(Y)}
    = \frac{\EV\bigl[\operatorname{Var}(Y \mid X_{\sim i})\bigr]}
    {\operatorname{Var}(Y)},
\end{equation}
where $X_{\sim i}$ denotes all inputs \emph{except} $X_i$.
The total-order index captures the variance attributable to $X_i$ including all
interactions with other inputs.
\end{definition}

\begin{keyidea}
$S_i$ measures the \emph{main effect} of input $i$: how much output variance
would be eliminated if $X_i$ were fixed at its true value.
$S_{T_i}$ measures the main effect \emph{plus all interactions}: fixing $X_i$
eliminates this fraction of output variance.
The difference $S_{T_i} - S_i$ quantifies the importance of interactions
involving $X_i$.
If all indices are close to their first-order values ($S_{T_i} \approx S_i$
for all $i$), the model is approximately additive and interactions are negligible.
\end{keyidea}

\subsection{Saltelli Sampling Design}

Efficient estimation of Sobol indices uses the Saltelli (2002) sampling design.
For $k$ inputs and $N$ base samples, this requires $N(k + 2)$ model evaluations:

\begin{algorithmbox}[title={Algorithm 2: Saltelli Sampling for Sobol Indices}]
\label{alg:saltelli}
\textbf{Input:} Number of base samples $N$, number of inputs $k$, input
distributions $\{F_1, \ldots, F_k\}$.
\begin{enumerate}
    \item Generate two independent $N \times k$ sample matrices $\bm{A}$ and
          $\bm{B}$, where each column is drawn from the corresponding input
          distribution.
    \item For each input $i = 1, \ldots, k$, construct matrix $\bm{A}_B^{(i)}$:
          take $\bm{A}$ and replace column $i$ with column $i$ from $\bm{B}$.
    \item Evaluate the model at all $N(k + 2)$ points:
          $f(\bm{A})$, $f(\bm{B})$, and $f(\bm{A}_B^{(i)})$ for each $i$.
    \item Estimate first-order indices:
          \begin{equation}
              \hat{S}_i = \frac{\frac{1}{N}\sum_{j=1}^{N}
              f(\bm{B})_j \bigl(f(\bm{A}_B^{(i)})_j - f(\bm{A})_j\bigr)}
              {\operatorname{Var}(Y)}.
          \end{equation}
    \item Estimate total-order indices:
          \begin{equation}
              \hat{S}_{T_i} = \frac{\frac{1}{2N}\sum_{j=1}^{N}
              \bigl(f(\bm{A})_j - f(\bm{A}_B^{(i)})_j\bigr)^2}
              {\operatorname{Var}(Y)}.
          \end{equation}
\end{enumerate}
\end{algorithmbox}

\subsection{Worked Example: Climate VaR Sensitivity}

\begin{example}[Sobol decomposition of Climate VaR]\label{ex:sobol}
We apply Sobol analysis to the stochastic Climate VaR model
(Algorithm~\ref{alg:climate-var}) with five uncertain inputs:
\begin{center}
\begin{tabular}{@{}lccccc@{}}
\toprule
\textbf{Input} & \textbf{Distribution} & $S_i$ & $S_{T_i}$ &
$S_{T_i} - S_i$ \\
\midrule
Carbon price ($p_{\text{CO}_2}$) & LogNormal($\mu$, $\sigma$) & 0.45 & 0.52 & 0.07 \\
Temperature ($T$)                & Normal($\mu_T$, $\sigma_T$) & 0.25 & 0.30 & 0.05 \\
GDP growth ($g$)                 & Normal($\mu_g$, $\sigma_g$) & 0.15 & 0.18 & 0.03 \\
Emission intensity ($e$)         & Uniform(0.8$e_0$, 1.2$e_0$) & 0.08 & 0.10 & 0.02 \\
Discount rate ($r$)              & Uniform(2\%, 5\%)            & 0.04 & 0.05 & 0.01 \\
\midrule
\textbf{Sum} & & 0.97 & 1.15 & \\
\bottomrule
\end{tabular}
\end{center}

\textbf{Interpretation:}
\begin{itemize}
    \item \textbf{Carbon price} is the dominant driver, accounting for 45\% of
          output variance on its own and 52\% including interactions.
          Policy uncertainty about the trajectory and level of carbon pricing
          dominates the Climate VaR uncertainty.
    \item \textbf{Temperature} is the second-largest driver (25\%), primarily
          affecting the physical risk channel through the quadratic damage function.
    \item \textbf{GDP growth} contributes 15\%, affecting both revenue projections
          and the denominator of the damage function.
    \item The sum of first-order indices ($\sum S_i = 0.97$) is close to 1,
          indicating that the model is approximately \emph{additive} ---
          interactions between inputs are modest.
    \item The sum of total-order indices exceeds 1 ($\sum S_{T_i} = 1.15$)
          because total-order indices double-count shared interaction variance.
\end{itemize}
\end{example}

\begin{keyresult}
Sobol analysis reveals that \textbf{carbon pricing policy} is the single most
important driver of climate portfolio risk --- more important than the physical
temperature outcome.
This finding has a direct implication: risk managers should focus scenario
analysis on the range of plausible carbon price trajectories and treat
physical risk as a secondary (but still material) driver.
For portfolios with high physical exposure (real estate, agriculture), the
ranking may reverse.
\end{keyresult}


%% ============================================================================
\section{TCFD Reporting Metrics}
\label{sec:tcfd}
%% ============================================================================

The TCFD framework (and its successor, ISSB IFRS~S2) defines specific metrics
for disclosing portfolio-level climate exposure.
These metrics complement the scenario analysis in Sections~\ref{sec:transition-risk}--\ref{sec:climate-var}
with standardized, comparable measures.

\subsection{Weighted Average Carbon Intensity (WACI)}

\begin{definition}[WACI]\label{def:waci}
The \textbf{Weighted Average Carbon Intensity} of a portfolio is:
\begin{equation}\label{eq:waci}
    \text{WACI} = \sum_{i=1}^{N} w_i \cdot \text{CI}_i,
\end{equation}
where $w_i$ is the portfolio weight of holding $i$ and $\text{CI}_i$ is the
carbon intensity of the issuer (tCO$_2$/\$M revenue, Scope 1+2).
\end{definition}

WACI is the TCFD's recommended cross-sector metric because it:
\begin{itemize}
    \item Normalizes by revenue, making firms of different sizes comparable.
    \item Uses portfolio weights, reflecting actual exposure.
    \item Is readily computable from publicly available emissions data.
\end{itemize}

\begin{example}[Portfolio WACI computation]\label{ex:waci}
\begin{center}
\begin{tabular}{@{}lccc@{}}
\toprule
\textbf{Holding} & \textbf{Weight ($w_i$)} & \textbf{CI$_i$ (tCO$_2$/\$M)} &
\textbf{$w_i \times$ CI$_i$} \\
\midrule
Oil Major A      & 15\% & 480  & 72.0 \\
Utility B        & 10\% & 310  & 31.0 \\
Steel Co C       & 10\% & 220  & 22.0 \\
Bank D           & 25\% & 45   & 11.3 \\
Tech E           & 20\% & 8    & 1.6 \\
Healthcare F     & 20\% & 15   & 3.0 \\
\midrule
\textbf{Portfolio WACI} & & & \textbf{140.9} \\
\bottomrule
\end{tabular}
\end{center}
The portfolio's WACI of 140.9 tCO$_2$/\$M revenue can be benchmarked against the
MSCI ACWI index ($\sim$130 tCO$_2$/\$M) or a Paris-aligned benchmark
($\sim$70 tCO$_2$/\$M).
\end{example}

\subsection{Financed Emissions}

While WACI measures \emph{intensity}, financed emissions measure \emph{absolute}
contribution to global emissions:

\begin{definition}[Financed emissions]\label{def:financed-emissions}
The \textbf{financed emissions} attributable to a portfolio are:
\begin{equation}\label{eq:financed-emissions}
    E_{\text{financed}} = \sum_{i=1}^{N}
    \underbrace{\frac{V_i}{\text{EVIC}_i}}_{\text{ownership share}}
    \times \underbrace{E_i}_{\text{company emissions}},
\end{equation}
where $V_i$ is the portfolio's investment in holding $i$, $\text{EVIC}_i$ is the
enterprise value including cash of the issuer, and $E_i$ is the issuer's total
Scope 1+2 emissions (tCO$_2$).
\end{definition}

\subsection{Emission Scopes}

Corporate greenhouse gas emissions are classified into three scopes under the
GHG Protocol:

\begin{table}[H]
\centering
\caption{GHG Protocol emission scopes.}
\label{tab:emission-scopes}
\begin{tabular}{@{}lp{7cm}l@{}}
\toprule
\textbf{Scope} & \textbf{Description} & \textbf{Data Quality} \\
\midrule
Scope 1 & Direct emissions from owned or controlled sources
          (e.g., combustion in facilities, company vehicles). & High \\[4pt]
Scope 2 & Indirect emissions from purchased electricity, steam,
          heating, and cooling. & High \\[4pt]
Scope 3 & All other indirect emissions in the value chain ---
          upstream (purchased goods, business travel) and downstream
          (use of sold products, end-of-life treatment). & Low (estimated) \\
\bottomrule
\end{tabular}
\end{table}

\begin{warningbox}
For many sectors, Scope~3 emissions dwarf Scope~1+2.
An oil company's Scope~1+2 covers extraction and refining; its Scope~3 includes
the combustion of all sold fuel --- typically 5--10$\times$ larger.
A bank's Scope~1+2 is negligible (office buildings), but its Scope~3 (financed
emissions from its loan book) is enormous.
Reporting only Scope~1+2 can drastically understate a portfolio's true climate
footprint.
ISSB IFRS~S2 requires Scope~3 disclosure ``to the extent that it is material''
--- in practice, this means most large firms must report it.
\end{warningbox}

\subsection{Portfolio Alignment: Implied Temperature Rise}

The most forward-looking TCFD metric is the portfolio's \textbf{implied
temperature rise} (ITR):

\begin{definition}[Implied Temperature Rise]\label{def:itr}
The ITR of a portfolio estimates the global warming that would result if the
entire economy had the same emission trajectory as the portfolio's holdings.
It is computed by:
\begin{enumerate}
    \item Projecting each holding's emissions through 2050 based on its
          disclosed targets and historical trajectory.
    \item Aggregating to a portfolio-level emission pathway.
    \item Mapping the pathway to a temperature outcome using a carbon budget
          model.
\end{enumerate}
\end{definition}

A portfolio aligned with the Paris Agreement should have an ITR $\leq$ 1.5$^\circ$C.
Portfolios with high fossil fuel exposure typically have ITR values of
2.5--3.5$^\circ$C.

\subsection{TCFD Reporting Template}

A complete TCFD-aligned climate risk report includes:

\begin{algorithmbox}[title={TCFD Disclosure Template}]
\label{alg:tcfd-template}
\begin{enumerate}
    \item \textbf{Governance:} Board oversight of climate risk; management's
          role in assessing and managing climate-related risks.
    \item \textbf{Strategy:} Climate-related risks and opportunities over
          short, medium, and long term; impact on business, strategy, and
          financial planning; resilience under different climate scenarios.
    \item \textbf{Risk Management:} Processes for identifying, assessing, and
          managing climate-related risks; integration into overall risk management.
    \item \textbf{Metrics and Targets:}
          \begin{itemize}
              \item WACI (Eq.~\ref{eq:waci}): \underline{\hspace{2cm}}
                    tCO$_2$/\$M revenue.
              \item Financed emissions (Eq.~\ref{eq:financed-emissions}):
                    \underline{\hspace{2cm}} tCO$_2$.
              \item Scope breakdown: Scope 1 \underline{\hspace{1cm}},
                    Scope 2 \underline{\hspace{1cm}},
                    Scope 3 \underline{\hspace{1cm}}.
              \item Climate VaR by scenario (Eq.~\ref{eq:climate-var}):
                    \begin{itemize}
                        \item Net Zero 2050: \underline{\hspace{1.5cm}}\%
                        \item Below 2$^\circ$C: \underline{\hspace{1.5cm}}\%
                        \item Delayed Transition: \underline{\hspace{1.5cm}}\%
                        \item Current Policies: \underline{\hspace{1.5cm}}\%
                    \end{itemize}
              \item Implied Temperature Rise: \underline{\hspace{1.5cm}}$^\circ$C.
              \item Reduction targets: \underline{\hspace{3cm}}.
          \end{itemize}
\end{enumerate}
\end{algorithmbox}

\begin{prereq}
This section connects directly to the code implementation.
The pipeline must:
(1) ingest firm-level emission data and carbon price paths from NGFS;
(2) compute WACI and financed emissions at the portfolio level;
(3) run the Climate VaR simulation (Algorithm~\ref{alg:climate-var}) for each
    scenario;
(4) perform Sobol sensitivity analysis (Algorithm~\ref{alg:saltelli});
(5) produce a TCFD-formatted output with all metrics populated.
\end{prereq}


%% ============================================================================
\section{References}
\label{sec:references}
%% ============================================================================

\begin{enumerate}[label={[\arabic*]}]
    \item \textbf{Task Force on Climate-related Financial Disclosures} (2017).
          \emph{Recommendations of the Task Force on Climate-related Financial
          Disclosures}.
          Financial Stability Board.

    \item \textbf{Network for Greening the Financial System} (2024).
          \emph{NGFS Climate Scenarios --- Phase~V}.
          NGFS Secretariat, Banque de France.

    \item \textbf{Nordhaus, W.~D.} (2017).
          Revisiting the social cost of carbon.
          \emph{Proceedings of the National Academy of Sciences},
          114(7):1518--1523.

    \item \textbf{Acharya, V.~V., Berner, R., and Engle, R.~F.} (2023).
          Climate Stress Testing.
          \emph{Annual Review of Financial Economics}, 15:291--326.

    \item \textbf{European Central Bank} (2025).
          Climate risk and financial stability: Insights from the 2025
          macroprudential stress test.
          \emph{ECB Macroprudential Bulletin}.

    \item \textbf{IPCC} (2021).
          \emph{Climate Change 2021: The Physical Science Basis}.
          Contribution of Working Group~I to the Sixth Assessment Report (AR6).
          Cambridge University Press.

    \item \textbf{International Sustainability Standards Board} (2023).
          \emph{IFRS S2: Climate-related Disclosures}.
          IFRS Foundation.

    \item \textbf{Howard, P.~H. and Sterner, T.} (2017).
          Few and not so far between: A meta-analysis of climate damage estimates.
          \emph{Environmental and Resource Economics}, 68(1):197--225.

    \item \textbf{Saltelli, A.} (2002).
          Making best use of model evaluations to compute sensitivity indices.
          \emph{Computer Physics Communications}, 145(2):280--297.

    \item \textbf{Sobol', I.~M.} (2001).
          Global sensitivity indices for nonlinear mathematical models and their
          Monte Carlo estimates.
          \emph{Mathematics and Computers in Simulation}, 55(1--3):271--280.
\end{enumerate}

\end{document}
